\documentclass[12pt]{article}
\usepackage[utf8]{inputenc}
\usepackage[T1]{fontenc}
\usepackage{hyperref}
\title{Wolaytta Language Resources for K-12 Curriculum Development}
\author{Compiled by Claude Sonnet 5 (claude-sonnet-5) for Andr\'as Kornai}
\date{}
\begin{document}
\maketitle

\section{Introduction}
Wolaytta (ISO 639-3 wal; Glottocode wola1242) is an Omotic language of Ethiopia with roughly 5.8 million speakers. It is used as a language of instruction. This search was brief and found reference and research material but no current textbook lists.

\section{Primary Educational Resources}
\begin{itemize}
\item Alemayehu Dogamo, ``The Challenges of Mother Tongue Education: The Case of Wolaytta'', Addis Ababa University master's thesis (2005). It covers primary schools in the Wolaytta Zone, teacher training, reading materials, terminology and orthography: \url{https://etd.aau.edu.et/items/8aebe3e6-dfd8-4ce8-ac10-b94cb33328cc}
\end{itemize}
No actual school textbooks or primers were located. Wikipedia reports that 1998 textbooks in the artificial ``Wegagoda'' were withdrawn after protest and replaced with Wolaytta materials.

\section{Linguistic and Reference Materials}
\begin{itemize}
\item Lamberti and Sottile, \emph{The Wolaytta Language} (1997, Studia Linguarum Africae Orientalis 6), a grammar with a Wolaytta--English lexicon of about 290 pages: \url{https://www.koeppe.de/titel_the-wolaytta-language}
\item Wikipedia entry on the Wolaytta language: \url{https://en.wikipedia.org/wiki/Wolaytta_language}
\item Wikipedia entry on the Wolaytta people: \url{https://en.wikipedia.org/wiki/Wolaytta_people}
\end{itemize}

\section{Religious and Community Resources}
Wikipedia reports a Wolaytta New Testament (1981) and a full Bible (2002). No direct scripture URL appeared in results.

\section{Digital and Audio Resources}
None found in this search.

\section{Curriculum Development Context}
Wolaytta has an established place in mother-tongue primary education, with documented difficulties around textbooks, terminology and teacher training. The thesis above is the best starting point. Current textbooks and digital resources still need to be located.
\end{document}